% !TeX root = main.tex
\IfFileExists{../apunte_mecanica_clasica.tex}{%
  \documentclass[../apunte_mecanica_clasica.tex]{subfiles}%
}{%
  \documentclass[apunte_mecanica_clasica.tex]{subfiles}%
}

\begin{document}

\chapter{Símbolos de Christoffel}


\section{Variación de las bases coordenadas}

En general, los vectores base coordenados 
$\{\vec{e}_i(x^j)\}_{i=1}^3$ son \textbf{campos vectoriales}, 
es decir, pueden variar punto a punto en el espacio, 
tanto en su dirección como en su magnitud.

Consideremos los vectores base coordenados 
$\vec{e}_i(x^j)$ y $\vec{e}_i(x^j + dx^j)$, evaluados en los 
puntos $P$ y $Q$, con coordenadas $x^i$ y $x^i + dx^i$, respectivamente.  

Podemos compararlos a primer orden en $dx^j$, para estudiar cómo 
cambian los vectores $\vec{e}_i(x^j)$ cuando pasamos de $x^j$ a 
$x^j + dx^j$.

\begin{figure}[h]
    \centering
    \includegraphics[width=0.4\linewidth]{figures/05_simbolos_de_christoffel/sim_fig1.png}
    \caption{Comparación entre los vectores base $\vec{e}_i$ en dos 
    puntos cercanos $P(x^j)$ y $Q(x^j + dx^j)$, mostrando cómo varían 
    con la posición en el espacio.}
\end{figure}

Luego, tenemos que el cambio es dado por (se usa la regla de la cadena):
\begin{align*}
    d\vec{e}_i &= \vec{e}_i(x^j + dx^j) - \vec{e}_i(x^j) \\[4pt]
               &= \frac{\partial \vec{e}_i}{\partial x^1}\, dx^1
                + \frac{\partial \vec{e}_i}{\partial x^2}\, dx^2
                + \frac{\partial \vec{e}_i}{\partial x^3}\, dx^3 \\[4pt]
               &= \sum_{j=1}^3 \frac{\partial \vec{e}_i}{\partial x^j}\, dx^j \,.
\end{align*}

El cambio de los vectores base coordenados es proporcional a
$\dfrac{\partial \vec{e}_i}{\partial x^j}$.  
Y como $d\vec{e}_i$ es un vector (diferencia de vectores), entonces 
$\dfrac{\partial \vec{e}_i}{\partial x^j}$ también es un vector.  
Por lo tanto, debe ser una combinación lineal de los vectores base 
coordenados.


\newpage

\section{Descomposición en la base coordenada}

Puesto que $\dfrac{\partial \vec{e}_i}{\partial x^j}$ posee dos 
índices libres ($i$ y $j$), entonces tiene $3 \times 3 = 9$ componentes.  

Para entender esto, consideremos el caso $i=1$:  
\begin{align*}
    \frac{\partial \vec{e}_1}{\partial x^1} 
        &= \alpha^1 \vec{e}_1 + \alpha^2 \vec{e}_2 + \alpha^3 \vec{e}_3
        = \sum_{k=1}^3 \alpha^k \vec{e}_k, \\[6pt]
    \frac{\partial \vec{e}_1}{\partial x^2} 
        &= \beta^1 \vec{e}_1 + \beta^2 \vec{e}_2 + \beta^3 \vec{e}_3
        = \sum_{k=1}^3 \beta^k \vec{e}_k, \\[6pt]
    \frac{\partial \vec{e}_1}{\partial x^3} 
        &= \gamma^1 \vec{e}_1 + \gamma^2 \vec{e}_2 + \gamma^3 \vec{e}_3
        = \sum_{k=1}^3 \gamma^k \vec{e}_k \,.
\end{align*}

De manera análoga, para $\vec{e}_2$ tenemos:
\begin{align*}
    \frac{\partial \vec{e}_2}{\partial x^1} 
        &= \lambda^1 \vec{e}_1 + \lambda^2 \vec{e}_2 + \lambda^3 \vec{e}_3
        = \sum_{k=1}^3 \lambda^k \vec{e}_k, \\[6pt]
    \frac{\partial \vec{e}_2}{\partial x^2} 
        &= \tau^1 \vec{e}_1 + \tau^2 \vec{e}_2 + \tau^3 \vec{e}_3
        = \sum_{k=1}^3 \tau^k \vec{e}_k, \\[6pt]
    \frac{\partial \vec{e}_2}{\partial x^3} 
        &= \epsilon^1 \vec{e}_1 + \epsilon^2 \vec{e}_2 + \epsilon^3 \vec{e}_3
        = \sum_{k=1}^3 \epsilon^k \vec{e}_k \,.
\end{align*}

Finalmente, para $\vec{e}_3$ se obtiene:
\begin{align*}
    \frac{\partial \vec{e}_3}{\partial x^1} 
        &= \pi^1 \vec{e}_1 + \pi^2 \vec{e}_2 + \pi^3 \vec{e}_3
        = \sum_{k=1}^3 \pi^k \vec{e}_k, \\[6pt]
    \frac{\partial \vec{e}_3}{\partial x^2} 
        &= \phi^1 \vec{e}_1 + \phi^2 \vec{e}_2 + \phi^3 \vec{e}_3
        = \sum_{k=1}^3 \phi^k \vec{e}_k, \\[6pt]
    \frac{\partial \vec{e}_3}{\partial x^3} 
        &= \psi^1 \vec{e}_1 + \psi^2 \vec{e}_2 + \psi^3 \vec{e}_3
        = \sum_{k=1}^3 \psi^k \vec{e}_k \,.
\end{align*}


\section{Notación matricial y definición}

Las expresiones anteriores se pueden escribir en una forma 
compacta y simple. La expresión $\frac{\partial \vec{e}_i}{\partial x^j}$   
se puede representar matricialmente, donde:
\begin{itemize}
    \item El índice $i$ rotula las \textbf{filas}.
    \item El índice $j$ rotula las \textbf{columnas}.
\end{itemize}

Entonces,
\begin{equation*}
    \frac{\partial \vec{e}_i}{\partial x^j} =
    \begin{pmatrix}
        \dfrac{\partial \vec{e}_1}{\partial x^1} & \dfrac{\partial \vec{e}_1}{\partial x^2} & \dfrac{\partial \vec{e}_1}{\partial x^3} \\[8pt]
        \dfrac{\partial \vec{e}_2}{\partial x^1} & \dfrac{\partial \vec{e}_2}{\partial x^2} & \dfrac{\partial \vec{e}_2}{\partial x^3} \\[8pt]
        \dfrac{\partial \vec{e}_3}{\partial x^1} & \dfrac{\partial \vec{e}_3}{\partial x^2} & \dfrac{\partial \vec{e}_3}{\partial x^3}
    \end{pmatrix}.
\end{equation*}

Sustituyendo las combinaciones lineales obtenidas:
\begin{equation*}
    \frac{\partial \vec{e}_i}{\partial x^j} =
    \begin{pmatrix}
        \sum_{k=1}^3 \alpha^k \vec{e}_k & \sum_{k=1}^3 \beta^k \vec{e}_k & \sum_{k=1}^3 \gamma^k \vec{e}_k \\[8pt]
        \sum_{k=1}^3 \lambda^k \vec{e}_k & \sum_{k=1}^3 \tau^k \vec{e}_k & \sum_{k=1}^3 \epsilon^k \vec{e}_k \\[8pt]
        \sum_{k=1}^3 \pi^k \vec{e}_k & \sum_{k=1}^3 \phi^k \vec{e}_k & \sum_{k=1}^3 \psi^k \vec{e}_k
    \end{pmatrix}.
\end{equation*}

En forma compacta:
\begin{equation*}
    \frac{\partial \vec{e}_i}{\partial x^j}
    = \sum_{k=1}^3 
    \begin{pmatrix}
        \alpha^k & \beta^k & \gamma^k \\
        \lambda^k & \tau^k & \epsilon^k \\
        \pi^k & \phi^k & \psi^k
    \end{pmatrix} \vec{e}_k \,.
\end{equation*}

Es claro que
\begin{equation*}
    \begin{pmatrix}
        \alpha^k & \beta^k & \gamma^k \\
        \lambda^k & \tau^k & \epsilon^k \\
        \pi^k & \phi^k & \psi^k
    \end{pmatrix}
\end{equation*}
es una matriz de $3 \times 3$ que denotaremos como
\begin{equation*}
    \Gamma^k{}_{ij} :=
    \begin{pmatrix}
        \Gamma^k{}_{11} & \Gamma^k{}_{12} & \Gamma^k{}_{13} \\
        \Gamma^k{}_{21} & \Gamma^k{}_{22} & \Gamma^k{}_{23} \\
        \Gamma^k{}_{31} & \Gamma^k{}_{32} & \Gamma^k{}_{33}
    \end{pmatrix}.
\end{equation*}

Por lo tanto, tenemos que
\begin{equation*}
    \frac{\partial \vec{e}_i}{\partial x^j} 
    = \sum_{k=1}^3 \Gamma^k{}_{ij} \, \vec{e}_k \,.
\end{equation*}

\begin{apuntebloque}{Definición}
Las cantidades $\Gamma^k{}_{ij}$ se llaman \textbf{símbolos de Christoffel}, 
y miden cómo cambian los vectores base coordenados a medida que son 
desplazados a lo largo de las líneas coordenadas.  

Los símbolos de Christoffel poseen tres índices, por lo tanto hay un total de
\[
3 \times 3 \times 3 = 27
\]
cantidades de símbolos de Christoffel.
\end{apuntebloque}


\section{Primeros ejemplos: cartesianas y esféricas}

\begin{apuntebloque}{Ejemplo: Símbolos de Christoffel en distintos sistemas coordenados}

\textbf{a) Coordenadas Cartesianas}  

En coordenadas cartesianas los vectores base 
$\{\hat{e}_i\}_{i=1}^3 = \{\hat{\imath}, \hat{\jmath}, \hat{k}\}$ son constantes.  
Por lo tanto:
\begin{equation*}
    \frac{\partial \vec{e}_i}{\partial x^j} = 0
    = \sum_{k=1}^3 \Gamma^k{}_{ij}\, \vec{e}_k \,.
\end{equation*}
Y como los $\{\hat{e}_i\}_{i=1}^3$ son linealmente independientes, se concluye que:
\begin{equation*}
    \Gamma^k{}_{ij} = 0, 
    \qquad i,j,k \in \{1,2,3\}.
\end{equation*}

\bigskip

\textbf{b) Coordenadas Esféricas}  

Tenemos que para $i=1, \,j=1$:
\begin{equation*}
   \frac{\partial \vec{e}_1}{\partial x^1} = \frac{\partial \vec{e}_r}{\partial r} = 0 
    = \Gamma^1{}_{11}\vec{e}_r + \Gamma^2{}_{11}\vec{e}_\phi + \Gamma^3{}_{11}\vec{e}_\theta \,,
\end{equation*}
y como los vectores base coordenados son linealmente independientes se tiene que
\begin{equation*}
\left.
    \begin{aligned}
        \Gamma^1{}_{11}&=0 \\
        \Gamma^2{}_{11}&=0 \\
        \Gamma^3{}_{11}&=0
    \end{aligned}
\right\}    
    \quad \Rightarrow \quad \Gamma^k{}_{11}=0, \; k=\{1,2,3\}.
\end{equation*}

Para el caso $i=1, \; j=2$:
\begin{equation*}
    \frac{\partial \vec{e}_1}{\partial x^2} =\frac{\partial \vec{e}_r}{\partial \phi} = \frac{1}{r}\vec{e}_\phi
    = \Gamma^1{}_{12}\vec{e}_r + \Gamma^2{}_{12}\vec{e}_\phi + \Gamma^3{}_{12}\vec{e}_\theta,
\end{equation*}
\begin{equation*}
    \therefore \quad 
    \left.
    \begin{aligned}
        &\Gamma^1{}_{12} =\Gamma^3{}_{12} = 0 \\
        &\Gamma^2{}_{12} = \tfrac{1}{r}
    \end{aligned}
    \right\}.
\end{equation*}

Para el caso $i=1, \; j=3$:
\begin{equation*}
    \frac{\partial \vec{e}_1}{\partial x^3} =\frac{\partial \vec{e}_r}{\partial \theta} = \frac{1}{r}\vec{e}_\theta
    = \Gamma^1{}_{13}\vec{e}_r + \Gamma^2{}_{13}\vec{e}_\phi + \Gamma^3{}_{13}\vec{e}_\theta,
\end{equation*}
\begin{equation*}
    \therefore \quad 
    \left.
    \begin{aligned}
        &\Gamma^1{}_{13} =\Gamma^2{}_{13} = 0 \\
        &\Gamma^3{}_{13} = \tfrac{1}{r}
    \end{aligned}
    \right\}.
\end{equation*}

Para el caso $i=2, \; j=1$:
\begin{equation*}
    \frac{\partial \vec{e}_2}{\partial x^1} =\frac{\partial \vec{e}_\phi}{\partial r} = \frac{1}{r}\vec{e}_\phi
    = \Gamma^1{}_{21}\vec{e}_r + \Gamma^2{}_{21}\vec{e}_\phi + \Gamma^3{}_{21}\vec{e}_\theta,
\end{equation*}
\begin{equation*}
    \therefore \quad 
    \left.
    \begin{aligned}
        &\Gamma^1{}_{21} =\Gamma^3{}_{21} = 0 \\
        &\Gamma^2{}_{21} = \tfrac{1}{r}
    \end{aligned}
    \right\}.
\end{equation*}

\end{apuntebloque}
\begin{apuntebloque}{}

Para el caso $i=2, \; j=2$:
\begin{equation*}
    \frac{\partial \vec{e}_2}{\partial x^2} =\frac{\partial \vec{e}_\phi}{\partial \phi} 
    = -r \sin^2\theta \,\vec{e}_r - \sin\theta\cos\theta \,\vec{e}_\theta=\Gamma^1{}_{22}\vec{e}_r+\Gamma^2{}_{22}\vec{e}_\phi+\Gamma^3{}_{22}\vec{e}_\theta,
\end{equation*}
\begin{equation*}
    \therefore \quad
    \left.
    \begin{aligned}
        &\Gamma^1{}_{22}=-r\sin^2\theta \\
        &\Gamma^2{}_{22}=0 \\ 
        &\Gamma^3{}_{22}=-\sin\theta\cos\theta
    \end{aligned}
    \right\}.
\end{equation*}

Para el caso $i=2, \; j=3$:
\begin{equation*}
    \frac{\partial \vec{e}_2}{\partial x^3} =\frac{\partial \vec{e}_\phi}{\partial \theta} = \cot\theta \, \vec{e}_\phi=\Gamma^1{}_{23}\vec{e}_r+\Gamma^2{}_{23}\vec{e}_\phi+\Gamma^3{}_{23}\vec{e}_\theta,
\end{equation*}
\begin{equation*}
    \therefore \quad 
    \left.
    \begin{aligned}
        &\Gamma^1{}_{23} =\Gamma^3{}_{23} = 0 \\
        &\Gamma^2{}_{23} = \cot\theta
    \end{aligned}
    \right\}.
\end{equation*}

Para el caso $i=3, \; j=1$:
\begin{equation*}
    \frac{\partial \vec{e}_3}{\partial x^1} =\frac{\partial \vec{e}_\theta}{\partial r} = \frac{1}{r}\vec{e}_\theta=\Gamma^1{}_{31}\vec{e}_r+\Gamma^2{}_{31}\vec{e}_\phi+\Gamma^3{}_{31}\vec{e}_\theta,
\end{equation*}
\begin{equation*}
    \therefore \quad 
    \left.
    \begin{aligned}
        &\Gamma^1{}_{31} =\Gamma^2{}_{31} = 0 \\
        &\Gamma^2{}_{31} =  \tfrac{1}{r}
    \end{aligned}
    \right\}.
\end{equation*}

Para el caso $i=3, \; j=2$:
\begin{equation*}
    \frac{\partial \vec{e}_3}{\partial x^2} =\frac{\partial \vec{e}_\theta}{\partial \phi} = \cot\theta \, \vec{e}_\phi=\Gamma^1{}_{32}\vec{e}_r+\Gamma^2{}_{32}\vec{e}_\phi+\Gamma^3{}_{32}\vec{e}_\theta,
\end{equation*}
\begin{equation*}
    \therefore \quad 
    \left.
    \begin{aligned}
        &\Gamma^1{}_{32} =\Gamma^3{}_{32} = 0 \\
        &\Gamma^2{}_{32} = \cot\theta
    \end{aligned}
    \right\}.
\end{equation*}

Para el caso $i=3, \; j=3$:
\begin{equation*}
    \frac{\partial \vec{e}_3}{\partial x^3} =\frac{\partial \vec{e}_\theta}{\partial \theta} = -r \vec{e}_r=\Gamma^1{}_{33}\vec{e}_r+\Gamma^2{}_{33}\vec{e}_\phi+\Gamma^3{}_{33}\vec{e}_\theta,
\end{equation*}
\begin{equation*}
    \therefore \quad 
    \left.
    \begin{aligned}
        &\Gamma^2{}_{33} =\Gamma^3{}_{33} = 0 \\
        &\Gamma^1{}_{33} = -r
    \end{aligned}
    \right\}.
\end{equation*}

\end{apuntebloque}


\section{Simetría en los índices inferiores}

En el cálculo de los símbolos de Christoffel se deben realizar 
$3 \times 3 \times 3 = 27$ cálculos, puesto que 
$\Gamma^k{}_{ij}$ se puede ver como una \textbf{matriz cúbica} 
de orden $3 \times 3 \times 3$, lo que puede volverse un cálculo largo.  

Sin embargo, no es necesario realizar todos los cálculos, 
puesto que los símbolos de Christoffel poseen una importante 
\textbf{propiedad de simetría}.

\begin{apuntebloque}{Proposición}
Los \textbf{símbolos de Christoffel} son simétricos en los índices inferiores:
\begin{equation*}
    \Gamma^k{}_{ij} = \Gamma^k{}_{ji}.
\end{equation*}

\textbf{Demostración.}  
De la definición de los símbolos de Christoffel se tiene que
\begin{align*}
    \sum_{k=1}^3 \Gamma^k{}_{ij}\,\vec{e}_k 
    &= \frac{\partial \vec{e}_i}{\partial x^j} \\
    &= \frac{\partial}{\partial x^j}\!\left(\frac{\partial \vec{x}}{\partial x^i}\right) \\
    &= \frac{\partial}{\partial x^i}\!\left(\frac{\partial \vec{x}}{\partial x^j}\right) \\
    &= \frac{\partial \vec{e}_j}{\partial x^i} \\
    &= \sum_{k=1}^3 \Gamma^k{}_{ji}\,\vec{e}_k \,.
\end{align*}

Se sigue  que,
\begin{equation*}
    \sum_{k=1}^3 \Gamma^k{}_{ij}\,\vec{e}_k - \sum_{k=1}^3 \Gamma^k{}_{ji}\,\vec{e}_k= 0.
\end{equation*}
Entonces, 
\begin{equation*}
    \sum_{k=1}^3 \big(\Gamma^k{}_{ij} - \Gamma^k{}_{ji}\big)\,\vec{e}_k = 0.
\end{equation*}

Como las bases coordenadas $\{\vec{e}_i\}_{i=1}^3$ son linealmente independientes, 
se concluye que
\begin{equation*}
    \Gamma^k{}_{ij} = \Gamma^k{}_{ji}.
\end{equation*}
\hfill $\square$
\end{apuntebloque}

La proposición anterior reduce el número de componentes 
independientes de $\Gamma^k{}_{ij}$, puesto que los índices 
$i$ y $j$ corresponden a una matriz simétrica.  

Luego, el número de componentes independientes en los 
índices $(i,j)$ es
\begin{equation*}
    \frac{1}{2}\, 3 \times (3+1) = 6.
\end{equation*}

Como el índice $k$ posee tres componentes, se concluye que 
$\Gamma^k{}_{ij}$ posee
\begin{equation*}
    6 \times 3 = 18
\end{equation*}
coeficientes independientes, y no $27$.

\begin{apuntebloque}{Ejemplo}

Del ejemplo de los coeficientes de Christoffel calculados para las 
coordenadas esféricas, y de acuerdo a la proposición demostrada, se concluye que:
\begin{align*}
    \Gamma^1{}_{12} = \Gamma^1{}_{21} &= 0, &
    \Gamma^2{}_{12} = \Gamma^2{}_{21} &= \tfrac{1}{r}, &
    \Gamma^3{}_{12} = \Gamma^3{}_{21} &= 0, \\[6pt]
    \Gamma^1{}_{13} = \Gamma^1{}_{31} &= 0, &
    \Gamma^2{}_{13} = \Gamma^2{}_{31} &= 0, &
    \Gamma^3{}_{13} = \Gamma^3{}_{31} &= \tfrac{1}{r}, \\[6pt]
    \Gamma^1{}_{23} = \Gamma^1{}_{32} &= 0, &
    \Gamma^2{}_{23} = \Gamma^2{}_{32} &= \cot\theta, &
    \Gamma^3{}_{23} = \Gamma^3{}_{32} &= 0.
\end{align*}


\end{apuntebloque}


\section{Fórmula en términos de la métrica}

Dada la base coordenada $\{\vec{e}_i\}_{i=1}^3$ se pueden obtener los 
coeficientes de Christoffel derivando la base con respecto a las 
coordenadas y luego usando la definición de los símbolos de Christoffel.  

Sin embargo, conocida la métrica y su inversa podemos determinar 
los símbolos de Christoffel usando la siguiente proposición.

\begin{apuntebloque}{Proposición}
Los \textbf{símbolos de Christoffel} satisfacen la propiedad:
\begin{equation*}
    \Gamma^k{}_{ij} = \tfrac{1}{2} \sum_{l=1}^3 g^{kl} 
    \left( \frac{\partial g_{il}}{\partial x^j} 
         + \frac{\partial g_{jl}}{\partial x^i} 
         - \frac{\partial g_{ij}}{\partial x^l} \right).
\end{equation*}

\textbf{Demostración.}  
Multiplicamos la definición de los símbolos de Christoffel por $\vec{e}_l$:
\begin{equation*}
    \frac{\partial \vec{e}_i}{\partial x^j} 
    = \sum_{k=1}^3 \Gamma^k{}_{ij}\,\vec{e}_k
    \quad \Longrightarrow \quad
    \frac{\partial \vec{e}_i}{\partial x^j}\cdot \vec{e}_l
    = \sum_{k=1}^3 \Gamma^k{}_{ij} (\vec{e}_k \cdot \vec{e}_l)
    = \sum_{k=1}^3 \Gamma^k{}_{ij}\, g_{kl}.
\end{equation*}

Ahora usamos que
\[
    g_{il} = \vec{e}_i \cdot \vec{e}_l,
\]
de modo que
\begin{align*}
    \frac{\partial g_{il}}{\partial x^j} 
    &= \frac{\partial}{\partial x^j}(\vec{e}_i \cdot \vec{e}_l) \\[4pt]
    &= \frac{\partial \vec{e}_i}{\partial x^j}\cdot \vec{e}_l 
       + \vec{e}_i \cdot \frac{\partial \vec{e}_l}{\partial x^j}.
\end{align*}

De aquí se deduce que
\begin{equation*}
    \frac{\partial \vec{e}_i}{\partial x^j}\cdot \vec{e}_l
    = \frac{\partial g_{il}}{\partial x^j} 
      - \vec{e}_i \cdot \frac{\partial \vec{e}_l}{\partial x^j}.
\end{equation*}

Aplicando nuevamente la definición de Christoffel en el segundo término:
\begin{equation*}
    \vec{e}_i \cdot \frac{\partial \vec{e}_l}{\partial x^j}
    = \sum_{k=1}^3 \Gamma^k{}_{lj} (\vec{e}_i \cdot \vec{e}_k)
    = \sum_{k=1}^3 \Gamma^k{}_{lj} g_{ik}.
\end{equation*}

Por lo tanto,
\begin{equation*}
    \frac{\partial g_{il}}{\partial x^j} 
    = \sum_{k=1}^3 \Gamma^k{}_{ij} g_{kl} 
      + \sum_{k=1}^3 \Gamma^k{}_{lj} g_{ik}.
\end{equation*}

Realizando permutaciones cíclicas en los índices $(i,j,l)$ se obtienen expresiones similares:
\begin{align*}
    \frac{\partial g_{jl}}{\partial x^i} 
    &= \sum_{k=1}^3 \Gamma^k{}_{ji} g_{kl} 
      + \sum_{k=1}^3 \Gamma^k{}_{li} g_{jk}, \\[6pt]
    \frac{\partial g_{ij}}{\partial x^l} 
    &= \sum_{k=1}^3 \Gamma^k{}_{il} g_{kj} 
      + \sum_{k=1}^3 \Gamma^k{}_{jl} g_{ik}.
\end{align*}

Combinando estas tres expresiones y usando la simetría de los símbolos de Christoffel ($\Gamma^k{}_{ij} = \Gamma^k{}_{ji}$) y la simetría del tensor métrico ($g_{ij} = g_{ji}$), se obtiene:
\begin{equation*}
    \frac{\partial g_{il}}{\partial x^j} 
    + \frac{\partial g_{jl}}{\partial x^i} 
    - \frac{\partial g_{ij}}{\partial x^l}
    = 2 \sum_{k=1}^3 \Gamma^k{}_{ij} g_{kl}.
\end{equation*}

Finalmente, contrayendo con $g^{nl}$ y sumando sobre $l$:
\begin{equation*}
    \sum_{l=1}^3 g^{nl} \left( \frac{\partial g_{il}}{\partial x^j} 
         + \frac{\partial g_{jl}}{\partial x^i} 
         - \frac{\partial g_{ij}}{\partial x^l} \right)
    = 2 \sum_{k=1}^3 \Gamma^k{}_{ij} \sum_{l=1}^3 g^{nl} g_{kl}
    = 2 \sum_{k=1}^3 \Gamma^k{}_{ij} \delta^n_k
    = 2 \Gamma^n{}_{ij}.
\end{equation*}

Por lo tanto,
\begin{equation*}
    \Gamma^n{}_{ij} = \tfrac{1}{2} \sum_{l=1}^3 g^{nl}
    \left( \frac{\partial g_{il}}{\partial x^j}
         + \frac{\partial g_{jl}}{\partial x^i}
         - \frac{\partial g_{ij}}{\partial x^l} \right).
\end{equation*}
Como $n$ es un índice libre, podemos renombrarlo como $k$ para obtener la expresión en la proposición.
\hfill $\square$
\end{apuntebloque}


\section{Contracción con la delta de Kronecker}

En la demostración anterior se usó la expresión:
\begin{equation}\label{eq:contraccion_delta}
    \sum_{k=1}^3 \Gamma^k{}_{ij}\,\delta^n{}_k = \Gamma^n{}_{ij},
\end{equation}
la cual no se había demostrado con anterioridad.  
Ahora procederemos a demostrarla.  

En la expresión \eqref{eq:contraccion_delta} la delta de Kronecker 
posee el índice $n$ libre, por lo que puede tomar los valores 
$1,2,3$.

\textbf{Demostración de la expresión \eqref{eq:contraccion_delta}.}  

En la expresión
\[
    \sum_{k=1}^3 \Gamma^k{}_{ij}\,\delta^n{}_k = \Gamma^n{}_{ij},
\]
el índice $n$ es libre, por lo que puede tomar los valores $1,2,3$:

\begin{itemize}
    \item[a)] Si $n=1$, entonces
    \begin{align*}
        \sum_{k=1}^3 \Gamma^k{}_{ij}\,\delta^1{}_k
        &= \Gamma^1{}_{ij}\,\delta^1{}_1 
         + \Gamma^2{}_{ij}\,\delta^1{}_2
         + \Gamma^3{}_{ij}\,\delta^1{}_3 \\[4pt]
        &= \Gamma^1{}_{ij}\cdot 1 + \Gamma^2{}_{ij}\cdot 0 + \Gamma^3{}_{ij}\cdot 0 \\[4pt]
        &= \Gamma^1{}_{ij}.
    \end{align*}

    \item[b)] Si $n=2$, entonces
    \begin{align*}
        \sum_{k=1}^3 \Gamma^k{}_{ij}\,\delta^2{}_k
        &= \Gamma^1{}_{ij}\,\delta^2{}_1 
         + \Gamma^2{}_{ij}\,\delta^2{}_2
         + \Gamma^3{}_{ij}\,\delta^2{}_3 \\[4pt]
        &= 0 + \Gamma^2{}_{ij}\cdot 1 + 0 \\[4pt]
        &= \Gamma^2{}_{ij}.
    \end{align*}

    \item[c)] Si $n=3$, entonces
    \begin{align*}
        \sum_{k=1}^3 \Gamma^k{}_{ij}\,\delta^3{}_k
        &= \Gamma^1{}_{ij}\,\delta^3{}_1 
         + \Gamma^2{}_{ij}\,\delta^3{}_2
         + \Gamma^3{}_{ij}\,\delta^3{}_3 \\[4pt]
        &= 0 + 0 + \Gamma^3{}_{ij}\cdot 1 \\[4pt]
        &= \Gamma^3{}_{ij}.
    \end{align*}
\end{itemize}

Resumiendo y compactando, se obtiene:
\begin{align*}
    \sum_{k=1}^3 \Gamma^k{}_{ij}\,\delta^1{}_k &= \Gamma^1{}_{ij}, \\[4pt]
    \sum_{k=1}^3 \Gamma^k{}_{ij}\,\delta^2{}_k &= \Gamma^2{}_{ij}, \\[4pt]
    \sum_{k=1}^3 \Gamma^k{}_{ij}\,\delta^3{}_k &= \Gamma^3{}_{ij}.
\end{align*}

Por lo tanto:
\begin{equation*}
    \sum_{k=1}^3 \Gamma^k{}_{ij}\,\delta^n{}_k = \Gamma^n{}_{ij},
    \qquad i,j,k \in \{1,2,3\}\,,
\end{equation*}
lo que demuestra la identidad \eqref{eq:contraccion_delta}. La expresión \eqref{eq:contraccion_delta} corresponde al producto de la 
matriz $\Gamma^k{}_{ij}$ por la matriz identidad $\delta^n{}_k$.

\bigskip

Cuando se usa la notación de índices, y no la matricial, 
la expresión \eqref{eq:contraccion_delta} se obtiene aplicando 
las siguientes dos reglas:

\begin{enumerate}
    \item La delta de Kronecker elimina el símbolo de suma, 
    que suma uno de los índices de la delta de Kronecker:
    \begin{equation}\label{eq:simbolos_1}
        \sum_{k=1}^3 \Gamma^k{}_{ij}\,\delta^n{}_k 
        = \Gamma^k{}_{ij}\,\delta^n{}_k.
    \end{equation}

    \item El índice repetido, bajo la suma, es reemplazado por el 
    índice libre presente en la delta de Kronecker. Además, 
    se elimina tanto el índice sumado como la delta de Kronecker.  
    De \eqref{eq:simbolos_1} se obtiene:
    \begin{equation*}
        \Gamma^k{}_{ij}\,\delta^n{}_k = \Gamma^n{}_{ij}.
    \end{equation*}
\end{enumerate}

Resumiendo las dos reglas se tiene
\begin{equation*}
    \sum_{k=1}^3 \Gamma^k{}_{ij}\,\delta^n{}_k = \Gamma^n{}_{ij}.
\end{equation*}

\bigskip

Estas dos reglas se aplican a todos los casos donde aparezca la 
\textbf{delta de Kronecker} (con índices abajo, arriba o mixtos) 
con uno o con sus dos índices sumados.  

\begin{enumerate}
    \item Para un vector:
    \begin{equation*}
        \sum_{j=1}^3 V^j \,\delta^i{}_j = V^i.
    \end{equation*}

    \item Para un tensor de tipo $(2,1)$:
    \begin{equation*}
        \sum_{l=1}^3 M^{ij}{}_l \,\delta^l{}_k = M^{ij}{}_k.
    \end{equation*}

    \item Para un tensor de tipo $(2,2)$:
    \begin{align*}
        \sum_{k=1}^3 \sum_{l=1}^3 T^{ki}{}_{jl} \,\delta^l{}_m \,\delta^n{}_k
        &= \sum_{k=1}^3 \left( \sum_{l=1}^3 T^{ki}{}_{jl} \,\delta^l{}_m \right)\,\delta^n{}_k \\[6pt]
        &= \sum_{k=1}^3 T^{ki}{}_{jm} \,\delta^n{}_k \\[6pt]
        &= T^{ni}{}_{jm}.
    \end{align*}
\end{enumerate}


\section{Ejemplos con la fórmula métrica}

Volviendo al teorema, sean los ejemplos:

\begin{apuntebloque}{Ejemplo}

\textbf{a) Coordenadas Cartesianas.}  
Para las coordenadas cartesianas la métrica $\delta_{ij}$ y su inversa 
$\delta^{ij}$ son constantes, entonces:
\begin{align*}
    \Gamma^n{}_{ij} &= \tfrac{1}{2} \sum_{l=1}^3 g^{nl} 
    \left( \frac{\partial \delta_{il}}{\partial x^j}
         + \frac{\partial \delta_{jl}}{\partial x^i}
         - \frac{\partial \delta_{ij}}{\partial x^l} \right) \\[6pt]
    &= \tfrac{1}{2} \sum_{l=1}^3 \delta^{nl} (0+0-0)\\
    &= 0.
\end{align*}
Por lo tanto, en coordenadas cartesianas se cumple:
\begin{equation*}
    \Gamma^n{}_{ij} = 0.
\end{equation*}

\end{apuntebloque}

\begin{apuntebloque}{}

\textbf{b) Coordenadas Esféricas.}  
En coordenadas esféricas, la métrica y su inversa son:
\begin{equation*}
    g_{ij} =
    \begin{pmatrix}
        1 & 0 & 0 \\
        0 & r^2 \sin^2\theta & 0 \\
        0 & 0 & r^2
    \end{pmatrix}\qquad\text{y}
    \qquad
    g^{ij} =
    \begin{pmatrix}
        1 & 0 & 0 \\
        0 & \tfrac{1}{r^2 \sin^2\theta} & 0 \\
        0 & 0 & \tfrac{1}{r^2}
    \end{pmatrix}.
\end{equation*}

\textit{Caso $i=1, j=1$:}
\begin{align*}
    \Gamma^k{}_{11} 
    &= \frac{1}{2} \sum_{l=1}^3 g^{kl}
    \left( \frac{\partial g_{1l}}{\partial x^1}
         + \frac{\partial g_{l1}}{\partial x^1}
         - \frac{\partial g_{11}}{\partial x^l} \right) \\[6pt]
    &= \frac{1}{2} \sum_{l=1}^3 g^{kl}
    \left( \frac{\partial g_{1l}}{\partial r}
         + \frac{\partial g_{l1}}{\partial r}
         - \frac{\partial g_{11}}{\partial x^l} \right) \\[6pt]
    &= \frac{1}{2} \sum_{l=1}^3 g^{kl}
    \left( 2\frac{\partial g_{1l}}{\partial r}- 0\right) \\[6pt]
    &= \sum_{l=1}^3 g^{kl} \frac{\partial g_{1l}}{\partial r}\\[6pt]
    &= g^{k1} \frac{\partial g_{11}}{\partial r}+g^{k2} \frac{\partial g_{12}}{\partial r}+g^{k3} \frac{\partial g_{13}}{\partial r}\\[6pt]
    &= g^{k1} \frac{\partial (1)}{\partial r}+g^{k2} \frac{\partial (0)}{\partial r}+g^{k3} \frac{\partial (0)}{\partial r}\\[6pt]
    &=0.
\end{align*}
Por lo tanto
\begin{equation*}
    \Gamma^k{}_{11} = 0.
\end{equation*}

\textit{Caso $i=1, j=2$:}
\begin{align*}
    \Gamma^k{}_{12} 
    &= \frac{1}{2} \sum_{l=1}^3 g^{kl}
    \left( \frac{\partial g_{1l}}{\partial x^2}
         + \frac{\partial g_{l2}}{\partial x^1}
         - \frac{\partial g_{12}}{\partial x^l} \right) \\[6pt]
    &= \frac{1}{2} \sum_{l=1}^3 g^{kl}
    \left( \frac{\partial g_{1l}}{\partial \phi}
         + \frac{\partial g_{l2}}{\partial r}
         - \frac{\partial (0)}{\partial x^l} \right) \\[6pt]
    &= \frac{1}{2} \sum_{l=1}^3 g^{kl}
    \left( \frac{\partial g_{1l}}{\partial \phi}
         + \frac{\partial g_{l2}}{\partial r}\right) \\[6pt]
    &= \frac{1}{2} \left[  g^{k1}\left(\frac{\partial g_{11}}{\partial \phi}
         + \frac{\partial g_{12}}{\partial r}\right) 
         + g^{k2}\left(\frac{\partial g_{12}}{\partial \phi}
         + \frac{\partial g_{22}}{\partial r}\right)
         + g^{k3}\left(\frac{\partial g_{13}}{\partial \phi}
         + \frac{\partial g_{32}}{\partial r}\right)\right] \\[6pt]
    &= \frac{1}{2} \left[  g^{k1}\left(\frac{\partial (r)}{\partial \phi}
         + \frac{\partial (0)}{\partial r}\right) 
         + g^{k2}\left(\frac{\partial (0)}{\partial \phi}
         + \frac{\partial (r^2\sin^2\theta)}{\partial r}\right)
         + g^{k3}\left(\frac{\partial (0)}{\partial \phi}
         + \frac{\partial (0)}{\partial r}\right)\right] \\[6pt]
    &= \frac{1}{2} \left(  g^{k1}\cdot0 
         + g^{k2}\cdot2r\sin^2\theta
         + g^{k3}\cdot0\right) \\[6pt]
         &= g^{k2}r\sin^2\theta .
\end{align*}
Se sigue que:
\begin{align*}
    \Gamma^1{}_{12} = \Gamma^1{}_{21} &= g^{12}\, r\sin^2\theta = 0, \\[6pt]
    \Gamma^2{}_{12} = \Gamma^2{}_{21} &= g^{22}\, r\sin^2\theta 
    = \frac{1}{r^2\sin^2\theta}\,(r\sin^2\theta) 
    = \frac{1}{r}, \\[6pt]
    \Gamma^3{}_{12} = \Gamma^3{}_{21} &= g^{32}\, r\sin^2\theta = 0.
\end{align*}

\medskip
De forma análoga se procede con los demás valores de los índices.

\end{apuntebloque}


\section{Contracción de símbolos de Christoffel}

La siguiente proposición es importante para la construcción de los 
operadores clásicos.

\begin{apuntebloque}{Proposición}
Los \textbf{símbolos de Christoffel} satisfacen la ecuación:
\begin{equation*}
    \sum_{j=1}^3 \Gamma^j{}_{ij} 
    = \frac{\partial}{\partial x^i} \ln \sqrt{g},
\end{equation*}
donde $g = \det(g_{ij})$.

\medskip

\end{apuntebloque}

Notar que el índice $j$ está sumado en los símbolos de Christoffel.

\begin{apuntebloque}{Ejemplos}

\textbf{a) Coordenadas Cartesianas.} Para las coordenadas cartesianas la métrica es 
$g_{ij} = \delta_{ij}$ y su determinante es $g = 1$. Entonces:
\begin{align*}
    \sum_{j=1}^3 \Gamma^j{}_{ij} 
    &= \frac{\partial}{\partial x^i} \ln \sqrt{g} \\[6pt]
    &= \frac{\partial}{\partial x^i} \ln(1) \\[6pt]
    &= 0.
\end{align*}

Este resultado no debe sorprender, puesto que en coordenadas 
cartesianas se cumple directamente que:
\[
    \Gamma^k{}_{ij} = 0.
\]

\textbf{b) Coordenadas Esféricas.}  Para las coordenadas esféricas, la métrica es
\[
    g_{ij} =
    \begin{pmatrix}
        1 & 0 & 0 \\
        0 & r^2 \sin^2\theta & 0 \\
        0 & 0 & r^2
    \end{pmatrix},
\]
y su determinante es
\[
    g = r^4 \sin^2\theta.
\]

Así,
\begin{align*}
    \sum_{j=1}^3 \Gamma^j{}_{ij}
    &= \frac{\partial}{\partial x^i} \ln \sqrt{g} \\
    &= \frac{\partial}{\partial x^i} \ln \left(\sqrt{r^4 \sin^2\theta} \right) \\[6pt]
    &= \frac{\partial}{\partial x^i} \ln \left(r^2 \sin\theta \right) \\
    &= 2 \frac{\partial}{\partial x^i} \ln r
       + \frac{\partial}{\partial x^i} \ln \left(\sin\theta\right).
\end{align*}

Por lo tanto:
\begin{equation*}
    \sum_{j=1}^3 \Gamma^j{}_{ij}=2 \frac{\partial}{\partial x^i} \ln r
       + \frac{\partial}{\partial x^i} \ln \left(\sin\theta\right).
\end{equation*}

Ahora calculemos de forma explícita los casos $i\in\{1,2,3\}$.

\end{apuntebloque}

\begin{apuntebloque}{}

\bigskip
\textbf{Caso $i=1$:}
\begin{align*}
    \sum_{j=1}^3 \Gamma^j{}_{1j}
    &= 2 \frac{\partial}{\partial x^1} \ln r 
       + \frac{\partial}{\partial x^1} \ln (\sin\theta) \\
    &= 2 \frac{\partial}{\partial r} \ln r 
       + \frac{\partial}{\partial r} \ln (\sin\theta) \\
    &= \frac{2}{r} + 0 \\[6pt]
    &= \frac{2}{r}.
\end{align*}
Calculemos de forma explícita la expresión usando los símbolos de Christoffel 
calculados para las coordenadas esféricas, en particular 
$\sum_{j=1}^3 \Gamma^j{}_{1j}$, y de forma análoga para los demás casos.

\begin{align*}
    \sum_{j=1}^3 \Gamma^j{}_{1j}
    &= \Gamma^1{}_{11} + \Gamma^2{}_{12} + \Gamma^3{}_{13} \\
    &= 0 + \frac{1}{r} + \frac{1}{r} \\
    &= \frac{2}{r}. \\
\end{align*}

\textbf{Caso $i=2$:}
\begin{align*}
    \sum_{j=1}^3 \Gamma^j{}_{2j}
    &= 2 \frac{\partial}{\partial x^2} \ln r 
       + \frac{\partial}{\partial x^2} \ln (\sin\theta) \\
    &= 2 \frac{\partial}{\partial \phi} \ln r 
       + \frac{\partial}{\partial \phi} \ln (\sin\theta) \\
    &= 0 + 0 \\[4pt]
    &= 0.
\end{align*}
Calculemos de forma explícita la expresión usando los símbolos de Christoffel 
calculados para las coordenadas esféricas, en particular 
$\sum_{j=1}^3 \Gamma^j{}_{2j}$, y de forma análoga para los demás casos.

\begin{align*}
    \sum_{j=1}^3 \Gamma^j{}_{2j}
    &= \Gamma^1{}_{21} + \Gamma^2{}_{22} + \Gamma^3{}_{23}\\
    &= 0 + 0 + 0 \\[4pt]
    &= 0.
\end{align*}

\textbf{Caso $i=3$:}
\begin{align*}
    \sum_{j=1}^3 \Gamma^j{}_{3j}
    &= 2 \frac{\partial}{\partial \theta} \ln r 
       + \frac{\partial}{\partial \theta} \ln (\sin\theta) \\[4pt]
    &= 0 + \cot\theta \\[4pt]
    &= \cot\theta.
\end{align*}
Calculemos de forma explícita la expresión usando los símbolos de Christoffel 
calculados para las coordenadas esféricas, en particular 
$\sum_{j=1}^3 \Gamma^j{}_{3j}$, y de forma análoga para los demás casos.

\[
    \sum_{j=1}^3 \Gamma^j{}_{3j}
    = \Gamma^1{}_{31} + \Gamma^2{}_{32} + \Gamma^3{}_{33}
    = 0 + \cot\theta + 0 = \cot\theta.
\]

\end{apuntebloque}


\newpage

\end{document}
